\begin{claim}{D.2.18}{Proved}{The exact degree-eight group and the seven-skeleton extension}
Retain the notation of D.2.17, and let
\[
 \overline\alpha:S^5\longrightarrow T=\Omega^2\Sigma^2B
\]
be the double adjunct of $u_4\nu_4:S^7\to\Sigma^2B$, where
$u_4:S^4\to\Sigma^2B$ is the bottom class. Then
\[
 \pi_5T=\ZZ/4\{\overline\alpha\},\qquad
 \pi_6T=\ZZ/4\{\rho\},\qquad
 2\rho=\overline\alpha\eta_5\ne0.
 \tag{D.2.18a}
\]
Here $\rho$ can be \emph{any} element whose mod-two Hurewicz image is
$r=b^2+\{a,b\}+p_3$, including the actual Pontryagin--Thom element
of D.2.17. In particular, no element with that Hurewicz image has
order two. Equivalently,
\[
 \pi_8(\Sigma^2K(\FF_2,2))=\ZZ/4\{\widehat\rho\},\qquad
 2\widehat\rho=u_4\nu_4\eta_7.
 \tag{D.2.18b}
\]

There is a CW model of $A=Q\Sigma B$ with finite seven-skeleton
$K_7$ for which the map in D.2.17 can be chosen to extend to an
ordinary pointed map $G_7:K_7\to T$. Its stated six-dimensional
homology values remain unchanged. This assertion concerns the
seven-skeleton; the eight-cell and later extension conditions are
not settled by it.
\end{claim}
\begin{proof}
Throughout the proof of the homotopy-group calculation we work
two-locally. The groups of $Y=\Sigma^2B$, its connected loop spaces,
and the connected fibers used below are finite two-primary in
positive degrees, by the Serre-class theorem. Thus the final
finite-group conclusions are also integral conclusions. The
low cohomology of $B$ and the free double-loop homology operations
are those already fixed in D.2.17; see
\cite[\S2, nos.~8--9]{I_Serre1953} and
\cite[Chapter~III, Theorems~1.2--1.4 and 3.1--3.2]{I_CohenLadaMay1976}.

\medskip\noindent
\textbf{A three-connected fiber.}
The bottom class defines a map $Y\to K(\FF_2,4)$. Its double
loop is an $H$-map
\[
 q:T\longrightarrow B.
\]
The double-suspension unit $j:B\to T$ is a section of $q$:
$qj$ is the identity on the unique nonzero homotopy group of $B$.
Put $F=\operatorname{fib}(q)$. Multiplication in the grouplike
$H$-space $T$ gives a space equivalence
\[
 B\times F\longrightarrow T,\qquad (b,x)\longmapsto j(b)x.
 \tag{D.2.18c}
\]
One obtains an inverse using $q$ and translation by the inverse
of $j(q(t))$. A multiplicative section is not required.
Freudenthal makes $j$ three-connected, and the fiber sequence
then makes $F$ three-connected. Consequently
\[
 \pi_iF\cong\pi_iT\quad(i\ge3),\qquad
 \pi_6F\cong\pi_6\Sigma^\infty F.
 \tag{D.2.18d}
\]
The second isomorphism places the required calculation in the
stable range, despite the original group's unstable description.

\medskip\noindent
\textbf{All fiber homology through degree seven.}
To avoid confusing the label $e\in H_5B$ with a six-dimensional
fiber class, set
\[
 A_4=a^2,\quad z_5=e+ab,\quad u_5=\xi a,\quad
 d_6=b^2,\quad e_6=b^2+\{a,b\},\quad f_6=p_3.
\]
The complete table for $\widetilde H_*(F;\FF_2)$ in these degrees is
\[
\begin{array}{c|c|c|c}
 \text{degree}&\text{basis}&\beta=\operatorname{Sq}_*^1
                    &\operatorname{Sq}_*^2\\ \hline
 4&A_4&0&0\\
 5&z_5,u_5&A_4,0&0,0\\
 6&d_6,e_6,f_6&0,0,0&A_4,A_4,A_4\\
 7&V=\xi b,\ R=\{a,c\},\ P,\ Q&e_6,0,f_6,0&u_5,0,z_5,z_5
\end{array}
 \tag{D.2.18e}
\]
Here, as in D.2.17,
\[
 P=g+ae+bc+a^2b,\qquad Q=h+ad,
\]
where $g,h$ are the labels dual to $u^2v,uw$ in $H^*B$.
These are fiber classes, rather than merely elements of the
linear kernel of $q_*$. Indeed the splitting
\textup{(D.2.18c)} is an isomorphism of homology coalgebras.
Primitives in the tensor product of two connected coalgebras
are sums of primitives in its two factors. All displayed
classes are primitive and map to zero in $H_*B$.
For $u_5$ this also follows from its spherical realization by
$\overline\alpha$ in D.2.17. The degree-six assertion follows
from the divided-power multiplication in $H_*B$; in degree
seven $H_*B$ has no primitives, since $H^*B$ has no
indecomposable generator in that degree.

For completeness, the full dimensions of $H_iT$, for
$i=4,5,6,7$, are $2,4,6,9$. The product splitting subtracts
the base and mixed terms and leaves fiber dimensions $1,2,3,4$.
Thus no degree-seven attachment is omitted by restricting to
the displayed primitive list. The square on $V$ is the Nishida
identity $\operatorname{Sq}_*^2(\xi b)=\xi a$. For the other
possibly nonzero entries, Cartan gives
\[
 \operatorname{Sq}_*^2g=\operatorname{Sq}_*^2h=e,
 \qquad
 \operatorname{Sq}_*^2(bc)=\operatorname{Sq}_*^2(ad)=ab,
\]
which verifies the two entries $z_5$. The remaining entries
are the coproduct, Cartan, and Bockstein calculations of D.2.17.
Its explicit integral configuration chain further gives
\[
 \beta_2[d_6]=[u_5]\ne0.
 \tag{D.2.18f}
\]
The universal coefficient theorem and these first and second
Bocksteins determine the integral groups:
\[
 H_4F=\ZZ/2,\qquad H_5F=\ZZ/4,\qquad
 H_6F=(\ZZ/2)^2.
 \tag{D.2.18g}
\]
In the last degree, $e_6,f_6$ are the two independent first
Bockstein images; $d_6$ is the Tor class belonging to the
order-four group in degree five.

\medskip\noindent
\textbf{A finite stable cell model, with every attachment detected.}
Take a finite cellular approximation to
$\Sigma^\infty F_{(2)}$ through dimension eight. Such an
approximation exists because the low homotopy groups are finitely
generated and $F$ is three-connected. In this finite range,
cancel each adjacent cell pair with odd cellular incidence.
An odd incidence is an equivalence of two-local sphere cells;
ordinary row and column elimination in the attaching cofiber
diagram splits off that contractible pair. The remaining
cellular differential is zero modulo two. Its cell counts
through dimension seven therefore equal the dimensions in
\textup{(D.2.18e)}. Dimension eight was included so that any
extraneous seven-cell killed by an odd eight-cell incidence is
also removed. Now discard cells of dimension eight. This does
not change $\pi_5$ or $\pi_6$. We denote the resulting stable
skeleta by $L^{(i)}$; no assertion about an unstable minimal
CW model is needed.

Using \textup{(D.2.18e)--(D.2.18g)}, normalize the integral
cellular differentials through dimension six as
\[
 \partial z_5=2A_4,\qquad \partial d_6=4u_5,
 \qquad \partial e_6=\partial f_6=0.
 \tag{D.2.18h}
\]
The changes can preserve every indicated mod-two label:
the nonzero second Bockstein fixes the odd coefficient of
$\partial d_6/4$, and any other multiple of $4u_5$ is eliminated
by an even multiple of $d_6$. Consequently
$L^{(5)}=P^5(2)\vee S^5_{u_5}$, where
$P^5(2)=M(\ZZ/2,4)$. Its stable fifth homotopy group is
$\ZZ/2\{i\eta_4\}\oplus\ZZ_{(2)}\{u_5\}$.
The square column of \textup{(D.2.18e)} therefore specifies
all three six-cell attachments: they have the same
$i\eta_4$ component, and the degree components are $4,0,0$.

Change six-cell basis to $e_6,d_6-e_6,r_6=f_6-e_6$.
This gives the stable wedge
\[
 L^{(6)}\simeq C\vee M(\ZZ/4,5)\vee S^6_{r_6},\qquad
 C=\operatorname{cofib}\bigl((2,\eta_4):S^4\vee S^5\to S^4\bigr).
 \tag{D.2.18i}
\]
The mod-two reduction of $r_6$ is exactly the spherical
class $r=e_6+f_6$.

Here is a direct computation of every possible seven-cell
attaching component. The first two stable stems are
$\pi_1^S=\ZZ/2\{\eta\}$ and
$\pi_2^S=\ZZ/2\{\eta^2\}$, with $\eta^2\ne0$
\cite[\S2, equation~(2.1), p.~149]{I_Oka1986}.
The cofiber sequence defining $C$ gives an injection
\[
 \pi_6C\longrightarrow
 \pi_6(C/S^4)=\ZZ/2\{\eta_5\}\oplus\ZZ_{(2)}\{\iota_6\}
\]
whose image is $\ZZ/2\oplus2\ZZ_{(2)}$.
Indeed the preceding map onto $\pi_6S^4$ is surjective
because it contains multiplication by $\eta$ on
$\pi_1^S$, and the following map is reduction modulo two
on the second summand. Hence
\[
 \pi_6C=\ZZ_{(2)}\{\gamma\}\oplus\ZZ/2\{\delta\},
 \qquad \gamma\longmapsto(0,2),\quad
 \delta\longmapsto(\eta_5,0).
 \tag{D.2.18j}
\]
In particular $h(\gamma)=2e_6$, while the square from
the five-cell $z_5$ into a new seven-cell detects precisely
its $\delta$ coefficient. This last assertion follows by
collapsing the bottom $S^4$: the quotient is an isomorphism
on the cohomology degrees five and seven used by that square.
Similarly the degree-four Moore cofiber sequence gives
\[
 \pi_6M(\ZZ/4,5)=\ZZ/2\{t=i\eta_5\};
\]
the square from $u_5$ into the new seven-cell detects $t$.
Finally $\pi_6S^6=\ZZ_{(2)}\{\rho_0\}$.
Thus the integral degrees and these two square coefficients
determine each attachment completely. There is no additional
secondary component left unspecified.

The degree-six cycles in the original cell coordinates are
freely generated by $e_6,f_6$. Their boundary subgroup is
$2\ZZ_{(2)}\{e_6,f_6\}$, by \textup{(D.2.18g)}.
After division by two the degree-seven boundary matrix is a
surjection onto a free rank-two module; its reduction modulo
two sends $V,P$ to $e_6,f_6$ and kills $Q,R$.
Choose preimages of the two basis vectors congruent to $V,P$
modulo two. An initial error is twice a vector, so surjectivity
corrects it by twice a chain. The kernel also reduces onto
the kernel modulo two, because the quotient is free. Hence
the zero-boundary basis can be chosen congruent to $Q,R$.
All square coefficients are preserved, and the four integral
boundaries are exactly $2e_6,2f_6,0,0$.

In the coordinates of \textup{(D.2.18i)}, the complete
attaching-map table is consequently
\[
\begin{array}{c|c}
 \text{seven-cell}&\text{attaching map in }\pi_6L^{(6)}\\ \hline
 V&\gamma+t\\
 P&\gamma+2\rho_0+\delta\\
 Q&\delta\\
 R&0.
\end{array}
 \tag{D.2.18k}
\]
Taking its cokernel gives
\[
 \pi_6F=
 \frac{\ZZ_{(2)}\gamma\oplus\ZZ/2\delta
       \oplus\ZZ/2t\oplus\ZZ_{(2)}\rho_0}
      {\langle\gamma+t,\ \gamma+2\rho_0+\delta,\ \delta\rangle}
 =\ZZ/4\{\rho_0\},\qquad 2\rho_0=t\ne0.
 \tag{D.2.18l}
\]
Also $\pi_5C=0$, so \textup{(D.2.18i)} gives
$\pi_5F=\ZZ/4$ with Hurewicz generator $u_5$; seven-cells
and higher cells do not change it. The class
$\overline\alpha$ has exactly this Hurewicz image by D.2.17,
and is therefore a generator. The image of $t$ is
$\overline\alpha\eta_5$; changing the odd generator of
$\pi_5F$ does not change this order-two element.
The Hurewicz image of $\rho_0$ is $r$. Every element with
that image is an odd generator of $\ZZ/4$, so its double
is the same nonzero $t$. This proves
\textup{(D.2.18a)--(D.2.18b)} for the actual geometric
representative as well as every other such representative.

\medskip\noindent
\textbf{The source quotient and the actual seven-stage extension.}
The map $C_2X\to A$ is seven-connected, as in D.2.17,
so choose a finite $K_7=A^{(7)}$ from a CW model of $C_2X$.
The canonical map $g_2=\Omega^2(\eta_2\wedge\id_B)$ is
already defined on it. The correction is zero on
$K_4=M(\ZZ/2,3)$, so its extension is a problem on the
four-connected quotient $Q_7=K_7/K_4$.

We record all relevant source cells. The six-skeleton of
the quotient has the Moore bottom $c_5=C$, Moore top
$d_6=D$, and free six-cycles
$e_6^{\rm src}=D+E$ and $z_6^{\rm src}=A_0^2$.
Here the integral lifts of $D,E$ have opposite boundaries
$4c_5,-4c_5$, respectively. The Snaith splitting gives
\[
 H_6(A;\ZZ)=H_5(B;\ZZ)\oplus
 H_6(D_2X;\ZZ)=(\ZZ/2)^2.
 \tag{D.2.18m}
\]
For the first term, the low dimensions from Serre and
$H_4(B;\ZZ)=\ZZ/4$ make $H_5B$ cyclic two-primary;
$\beta x=d+e\ne0$ fixes its exponent as two. The
order-four input is the Pontryagin-square calculation in
\cite[\S2.3 and proof of Lemma~6.1]{I_Mahowald1964}.
The weight-two term is its bottom group
$H_0(\Sigma_2;H_3X\otimes H_3X)=\ZZ/2$; there is no
lower term to support a differential. Higher weights begin
in degree nine. These uses of the stable splitting are
exactly \cite[Theorem~1.1 and \S2]{Snaith1974}.

The full degree-seven mod-two source basis is
\[
 F_7=\sigma f,\quad X_7=\sigma x,\quad
 R_7=Q^4A_0,\quad A_0B_0.
\]
Its Bocksteins are $0,e_6^{\rm src},z_6^{\rm src},
z_6^{\rm src}$, respectively. The product $A_0B_0$
must be retained here, even though it is not primitive.
On the unique degree-five cohomology class of the quotient,
\[
 \operatorname{Sq}^2(C^\vee)=\sigma(v^2).
 \tag{D.2.18n}
\]
It evaluates to one on $X_7$ and zero on the other three
basis elements. This follows on unit classes from
$\operatorname{Sq}^2(u^2)=v^2$, on $R_7$ from the
Nishida identity, and on $A_0B_0$ from Cartan.

Integral Smith reduction cancels any unit six/seven pairs.
Equation \textup{(D.2.18m)} and the Bocksteins normalize the
two remaining nonzero boundaries as
$2e_6^{\rm src},2z_6^{\rm src}$, preserving the specified
mod-two classes $X_7,R_7$. The zero-boundary classes include
$F_7$ and $A_0B_0+R_7$. Any additional seven-cell class
killed by higher cells has zero value under
\textup{(D.2.18n)}, by naturality.
There is no extra attaching invariant: the quotient is
four-connected and
\[
 \pi_6\bigl(M(\ZZ/4,5)\vee S^6_e\vee S^6_z\bigr)
 =\ZZ/2\{i\eta_5\}\oplus\ZZ\{e\}\oplus\ZZ\{z\}.
\]
The degree-five square detects the first summand exactly.
Thus the actual quotient has the homotopy type
\[
 \Bigl(M(\ZZ/4,5)\vee S^6_e\vee S^6_z\Bigr)
 \mathop{\cup}_{i\eta_5+2e}e^7
 \mathop{\cup}_{2z}e^7
 \ \vee\bigvee S^7.
 \tag{D.2.18o}
\]
In this dimension and connectivity the stable attachment
calculation is the unstable one; all possible Whitehead
products have larger dimension. Extra zero-degree
attachments vanish because their only possible $\eta$
component has just been checked to be zero.

Use the Moore extension
$\mu:M(\ZZ/4,5)\to T$ chosen in D.2.17, with bottom
$\overline\alpha$ and top mod-two image $k=\{a,b\}$.
On the two free six-spheres use $\rho$ and zero,
respectively. The only potentially nonzero obstruction
in \textup{(D.2.18o)} is
\[
 \overline\alpha\eta_5+2\rho=0,
 \tag{D.2.18p}
\]
by \textup{(D.2.18a)}. The other displayed attachment and
all free seven-sphere attachments impose no condition.
Choose nullhomotopies to extend the correction across the
seven-cells, pull it back along $K_7\to Q_7$, and add it
to $g_2$ using the multiplication of $T$.
On $C,D,E,A_0^2$ the correction has values
$\xi a,k,k+r,0$, respectively. Adding the canonical
values of D.2.17 gives precisely its displayed values,
including $D\mapsto b^2+k$ and $E\mapsto r$.
This constructs $G_7$ and finishes the stated finite-stage
extension.
\end{proof}
